\documentclass[conference]{IEEEtran}

\begin{document}

\title{Group 7 Project Abstract}

\author{
    \IEEEauthorblockN{Mohammad Islam, Macy Narang, Ken Jiang, and Jinxin Lin}
    \IEEEauthorblockA{
        \textit{Computer Science} \\
        \textit{CUNY Hunter College} \\
        mohammad.islam47@login.cuny.edu,\quad macy.narang99@login.cuny.edu \\
        ken.jiang24@login.cuny.edu,\quad jinxin.lin07@login.cuny.edu
    }
}

\maketitle

% Coarsely equalize the final columns as described in IEEEtran_HOWTO.
% Recheck this adjustment if the text changes.
\enlargethispage{-12.5\baselineskip}

\section{Introduction}
Several methods for estimating the win probability of a specific NFL team given
only the current state of an ongoing game are already widely used in sports
analytics. However, such calculations of win probability are not informative of
how that probability could change across the range of possible next plays and
their outcomes. Thus, the aim of this project is to address this problem by
estimating a distribution of win probabilities conditional on plausible next
plays. This service will be delivered through a web tool that can be used by
general audiences, professional analysts, and NFL personnel alike.

\section{Proposed Approach}
The models that calculate the probabilities mentioned above will be trained on
game-state data---including down, distance, time remaining, etc.---obtained from
datasets of historical NFL play-by-play data up to and including this season.
The models may be built using neural networks or empirical methods, and we will
evaluate their performance. The final model(s) should take vectors of live
game-state information obtained through web scraping as input, and deliver
estimated win probabilities corresponding to plausible play types as output.
The model(s) will be integrated into a web tool that allows users to make inputs
to the model by selecting an ongoing game and importing that game's live state,
or by manually entering information in fields that correspond to the model's
input variables. Upon receiving user
input, the webpage will display the resulting win probabilities associated with
each model.

\section{Intended Outcome}
If the models and web tool work as intended, they will give users insight into
how win probability changes throughout a game depending on what plays are called.
Thus, this tool can be used to assist football coaches in their play-making
decisions, for educational purposes among fans and new coaching personnel, and
to satisfy general curiosity among users. In certain situations, it could even
give people who participate in prediction markets an edge. If a trader has
already made an input to the web tool before a play has been completed, and if they
can then recognize what type of play is unfolding in real time, they would be
able to change how they are trading in response to the win probability
displayed on our website that is associated with that type of play. To further
aid users in usage and understanding of our win probability estimates, we will
display information and statistics about the players on the offensive team
that are relevant to the most successful play types, recent head-to-head matchup
data between the teams, and team rankings and statistics, among other things.
This data can be obtained using datasets, APIs, or web scraping.

\section{Conclusion}
Regardless of the setting, whether it be trading, education, or coaching, the
models' estimates should not override informed decision-making or replace
cognitive reasoning, but instead assist in them. Thus, the way in which the
tool is presented and curated should encourage users to think about how the
win probabilities should be considered within the context of the game state
they correspond to. This philosophy will inform our approach to designing the
website. If this approach is successful, our tool should be ready for
responsible and informed usage for multiple purposes by a wide range of people
among the millions that are interested in American football.

\end{document}
